\section{Weighted palettes and private resources}\label{sec:palettes}
\subsection{Compatibility}
Throughout this section, $Q$ is a finite simple graph with positive vertex weights $(x_v)_{v\in V(Q)}$ satisfying $\sum_vx_v=1$. For $U\subseteq V(Q)$ and $A\subseteq E(Q)$, write
\[
 x(U)=\sum_{v\in U}x_v,\qquad c(uv)=x_ux_v,\qquad
 c(A)=\sum_{e\in A}c(e).
\]
Put $D(v)=x(N(v))$, $\delta=\min_vD(v)$, and $m=c(E(Q))$. Neighborhoods are open, and walks and neighborhoods are always taken in $Q$.

\begin{definition}\label{def:compatibility}
Two distinct edge types $ab,cd\in E(Q)$ are \emph{compatible} if
\begin{equation}\label{eq:compatibility}
 (A_Q^2)_{ac}(A_Q^3)_{bd}=0
\end{equation}
for every choice of orientations $(a,b)$ and $(c,d)$ of the two edges. Here $A_Q$ is the adjacency matrix, so its powers count walks; walks may repeat vertices and may be closed. A \emph{pattern} is a nonempty pairwise compatible set of edge types.
\end{definition}
Compatibility is defined through walks rather than simple cycles. Section~\ref{sec:source} shows that, after the cleanup of Section~\ref{sec:cleanup}, the original color classes give patterns in the retained graph.

Let $T$ be the set of edges of $Q$ that lie in a triangle, and put $F=E(Q)\setminus T$. Let $Z$ be the set of vertices lying in no triangle, and define
\begin{equation}\label{eq:sectors}
 S=E(Q[Z]),\qquad J=F\setminus S,\qquad f=c(F).
\end{equation}
Thus $S\subseteq F$, and $f$ includes the mass of $S$.
For $A\subseteq E(Q)$, an \emph{exact fractional pattern cover} of $A$ is a collection of amounts $a_p\ge0$, indexed by patterns contained in $A$, with
\begin{equation}\label{eq:exactcover}
 \sum_{p\ni e}a_p=c(e)\quad(e\in A).
\end{equation}
Its cost is $\sum_pa_p$. Write $\Phi$ for the minimum cost of a cover of $E(Q)\setminus S$ and $\Phi_F$ for the minimum cost of a cover of $J$. The cover by singletons is feasible and there are finitely many patterns, so both minima are attained. Compatibility is always computed in $Q$, even when the covered set is smaller.

Lower bounds on covers come from weak duality. If $y_e\ge0$ and
\begin{equation}\label{eq:wholepattern}
 \sum_{e\in p}y_e\le1\quad\text{for every pattern }p\subseteq A,
\end{equation}
then every exact cover of $A$ has cost at least $\sum_{e\in A}c(e)y_e$. To see this, multiply~\eqref{eq:exactcover} by $y_e$ and interchange the two finite sums.

\subsection{Neighborhood geometry}
Both sectors are controlled by the following lemma.
\begin{lemma}\label{lem:cross}
For compatible edges $uv$ and $ab$, the nonempty intersections between the two indexed pairs
\[
 \{N(u),N(v)\}\quad\text{and}\quad\{N(a),N(b)\}
\]
form a partial matching: each member of either pair meets at most one member of the other. If this matching is perfect and the edges are oriented so that $N(u)\cap N(a)$ and $N(v)\cap N(b)$ are nonempty, then the crossed intersections are empty and the aligned pairs $N(u),N(a)$ and $N(v),N(b)$ are anticomplete.
\end{lemma}
\begin{proof}
If $N(u)$ meets both $N(a)$ and $N(b)$, there is a two-walk from $u$ to $a$, and for $z\in N(u)\cap N(b)$ the walk $vuzb$ has length three. These complementary walks are forbidden by~\eqref{eq:compatibility}. Relabeling the endpoints and exchanging the two edges gives the remaining row and column restrictions. In the perfect case, a two-walk from $u$ to $a$ excludes every three-walk from $v$ to $b$; equivalently, there is no edge between $N(v)$ and $N(b)$. The other aligned pair is treated in the same way.
\end{proof}
Neighborhoods are indexed by endpoints, so they are counted separately even when two endpoints coincide. We shall use repeatedly that when $\delta>1/3$ no three vertex neighborhoods are pairwise disjoint, since their weights would sum to more than one.

\begin{lemma}[Sector separation]\label{lem:sectors}
Suppose $\delta>1/3$. No compatible pair mixes $T$ and $F$, and every compatible triple consists entirely of edges of $F$.
\end{lemma}
\begin{proof}
Suppose $uv$ is compatible with $ab\in F$. Since $N(a)$ and $N(b)$ are disjoint, neither $N(u)$ nor $N(v)$ can avoid both, for that would give three disjoint neighborhoods. By Lemma~\ref{lem:cross} the intersections therefore form a perfect matching; align the pairs as in that lemma. If $uv\in T$, choose $z\in N(u)\cap N(v)$. By anticompleteness $N(z)$ avoids both $N(a)$ and $N(b)$, again giving three disjoint neighborhoods. This proves the first assertion.

For the second, consider the six indexed endpoint neighborhoods of three compatible edges, grouped into three pairs. Between any two pairs the intersection relation is a partial matching. There is no independent transversal of the three pairs, because three selected neighborhoods with no pairwise intersections would be disjoint.

We claim that every cross matching is perfect. Suppose a member $A$ of the first pair meets neither member of the second. For each member $B$ of the second pair, every member of the third must meet $A$ or $B$, since otherwise these three form an independent transversal. Each of $A,B$ has at most one cross neighbor in the third pair. Hence $A$ meets one member of the third pair, and \emph{both} members of the second pair must meet the other. This contradicts the partial-matching property.

Align the second and third pairs with the first, and denote their members by $A_i,B_i$ for $i=1,2,3$. The matching between pairs two and three is also aligned, since a crossed matching would leave $A_1,B_2,B_3$ pairwise disjoint. Thus, for distinct $i,j$, the sets $A_i,B_j$ are disjoint, while $A_i,A_j$ and $B_i,B_j$ are anticomplete by Lemma~\ref{lem:cross}. If the $i$th edge is triangular, take $z\in A_i\cap B_i$. For the two other indices $j,k$, the three neighborhoods $N(z),A_j,B_k$ are disjoint, a contradiction. Consequently all three edges lie in $F$.
\end{proof}

Let $\mathcal C_T$ be the graph whose vertices are the edge types in $T$ and whose edges are the compatible pairs. A \emph{capacitated fractional matching} in $\mathcal C_T$ assigns $z_{eg}\ge0$ to each unordered compatible pair $\{e,g\}$, subject to
\[
 \sum_{g:\{e,g\}\in E(\mathcal C_T)}z_{eg}\le c(e)\quad(e\in T).
\]
Let $\nu$ be the maximum of $\sum_{\{e,g\}}z_{eg}$, the value of a bounded linear program in finitely many variables.

\begin{proposition}[Exact palette decomposition]\label{prop:decomposition}
If $\delta>1/3$, then
\begin{equation}\label{eq:decomposition}
 \Phi=\Phi_F+c(T)-\nu=\Phi_F+m-f-\nu.
\end{equation}
\end{proposition}
\begin{proof}
By Lemma~\ref{lem:sectors}, each pattern lies wholly in $F$ or wholly in $T$, and a $T$-pattern is a singleton or a pair. The pair amounts of any exact $T$-cover form a feasible fractional matching, and if their sum is $v$, counting the covered capacity shows that the cost is $c(T)-v$. Conversely, a feasible matching extends to an exact cover of the same cost by covering its unused capacities with singletons. Minimizing gives $c(T)-\nu$, independently of the $J$-sector.
\end{proof}

\subsection{Private resources}
For $v\in V(Q)$ write $N_T(v)=\{u:uv\in T\}$ and $d_F(v)=x(N_F(v))$. Define
\begin{equation}\label{eq:resources}
 R_{uv}=N_T(u)\cup N_T(v)\quad(uv\in F),\qquad
 q(w)=\sum_{v:wv\in T}x_vd_F(v).
\end{equation}

\begin{lemma}[Private-resource prices]\label{lem:resources}
Suppose $\delta>1/3$. If $uv,ab\in F$ are compatible, then
\[
 R_{uv}\cap\bigl(N(a)\cup N(b)\bigr)=\varnothing.
\]
Consequently, for every $w$, the set
\[
 D_w=\{e\in J:w\in R_e\}
\]
meets every pattern in $J$ at most once, and
\begin{equation}\label{eq:resourcebound}
 c(D_w)=q(w),\qquad \Phi_F\ge\max_wq(w).
\end{equation}
\end{lemma}
\begin{proof}
The proof of Lemma~\ref{lem:sectors} gives a perfect alignment for the two edges. In that alignment $N(u)$ avoids $N(b)$, so $N_T(u)$ avoids $N(b)$. If $z\in N_T(u)\cap N(a)$, the triangular edge $uz$ has a common neighbor $t\in N(u)$ adjacent to $z$, and the edge $tz$ contradicts anticompleteness of $N(u),N(a)$. Similarly, $N_T(v)$ avoids both neighborhoods of the mate. This proves the asserted disjointness, and in particular $R_{uv}\cap R_{ab}=\varnothing$.

Hence a fixed $w$ lies in the resource of at most one edge of any pattern, and the indicator of $D_w$ satisfies~\eqref{eq:wholepattern}. Edges in $S$ have empty resources. Moreover, $N(u)\cap N(v)=\varnothing$ for $uv\in F$, so the two possible resource incidences of such an edge cannot both occur. Expanding the objective therefore gives
\[
 \sum_{uv\in F}x_ux_v\ind{w\in R_{uv}}
 =\sum_{v:wv\in T}x_v\sum_{u:uv\in F}x_u=q(w).
\]
The bound on $\Phi_F$ now follows from~\eqref{eq:wholepattern}.
\end{proof}
